\documentclass[german,10pt]{book}      \input{preambel}         \begin{document}  \setcounter{chapter}{3}

\chapter{Die Jahre 1925 und 1926}
% Kap x
\label{chap_1925}

\info{Thomas Filk}{5.10.2025}

Im Sommer 1925 erschien die ber\"uhmte \glqq Umdeutungsarbeit\grqq\ von 
Werner Heisenberg, in welcher er den Grundstein f\"ur die sogenannte
Matrizenmechanik legte \cite{Heisenberg}. Noch im selben Jahr folgten zwei weitere Arbeiten, einmal
von Max Born und Pascual Jordan \cite{BornJordan}, und dann die 
bekannte \glqq Drei M\"anner Arbeit\grqq\ von Born, Jordan
und Heisenberg \cite{DreiMann}. In diesen Arbeiten wurde die Matrizenmechanik
nahezu in ihrer heutigen Form formuliert. Insbesondere findet man schon in der
Arbeit von Born und Jordan die kanonischen Vertauschungsrelationen f\"ur die
\glqq Matrizen\grqq\ zum Orts- und Impulsoperator, die in der fr\"uheren
Arbeit von Heisenberg nur andeutungsweise auftauchen. Noch im November
1925 schrieb Paul Dirac einen Artikel zu den \glqq Grundgleichungen der Quantentheorie\grqq, 
in dem er sich auf Heisenberg bezog. Im Januar 1926 konnte Wolfgang
Pauli das Spektrum des Wasserstoffatoms aus der
Matrizenalgebra ableiten. 

Anfang 1926 erschien dann die erste von insgesamt
vier \glqq Mitteilungen\grqq\ von Erwin Schr\"odinger mit dem Titel \glqq Quantisierung
als Eigenwertproblem\grqq, in welcher Schr\"odinger seine Wellenmechanik 
formulierte \cite{Schroedinger1bis4}. Schon sehr bald (April 1926) erschien eine weitere Arbeit von 
Schr\"odinger, in der er die \"Aquivalenz zwischen dem Formalismus von Heisenberg
und seiner Wellenmechanik zeigen konnte \cite{Schroedinger26}. 

Auch wenn wir heute die Quantenmechanik in einer ausgereifteren Form
lehren und verwenden, lohnt es sich, einen Blick in die Originalarbeiten zu werfen und
die Gedankeng\"ange von Heisenberg und Schr\"odinger nachzuvollziehen. 

\section{Die Matrizenmechanik von 1925}

Schon das sehr kurze Abstract der Heisenberg'schen Umdeutungsarbeit deutet die Absicht 
Heisenbergs an, eine neue Theorie der Mechanik ausschlie\ss lich auf beobachtbare
Gr\"o\ss en aufzubauen:
\vspace{0.3cm}

\noindent
\textit{In der Arbeit soll versucht werden, Grundlagen zu gewinnen f\"ur eine quantentheoretische
Mechanik, die ausschlie\ss lich auf Beziehungen zwischen prinzipiell beobachtbaren Gr\"o\ss en
basiert.} 
\vspace{0.3cm}

\noindent
Die Gr\"unde f\"ur diesen Ansatz lagen in der \"alteren Atomphysik, auf die wir kurz eingehen.

\subsection{Die Bohr-Sommerfeld'schen Atommodelle}

1913 hatte Niels Bohr ein Modell f\"ur das Wasserstoffatom vorgeschlagen, bei dem
sich ein Elektron auf einer Kreisbahn um einen positiv geladenen Atomkern bewegt \cite{Bohr1}.
Dies entsprach auch dem zwei Jahre zuvor entwickelten 
Atommodell von Rutherford. Allerdings hatte Bohr eine einschr\"ankende
Bedingung an die Kreisbahnen postuliert, in der die Konstante $h$, das Planck'sche
Wirkungsquantum, auftauchte. Insofern kann man das Bohr'sche Atommodell als
erstes Quantenmodell des Atoms auffassen. 

F\"ur Kreisbahnen ist die Bohr'sche Quantisierungsbedingung \"aquivalent zu
folgender Bedingung an den Drehimpuls $L$ des Elektrons:
\begin{equation}
              L = r_n \cdot p_n = n \hbar \, ,
\end{equation} 
wobei $r_n$ und $p_n$ der Radius und der zugeh\"orige Impuls der Kreisbahn zu
einer Quantenzahl $n$ sind. Erlaubte Kreisbahnen sind somit solche, bei denen der
Bahndrehimpuls ein ganzzahliges Vielfaches von $\hbar$ ist.
Zusammen mit der klassischen Bedingung, dass bei einer Kreisbahn die
Impuls\"anderung (manchmal spricht man auch von der Fliehkraft) $\dot{p}=mv^2/r=p^2/mr$ gleich der
Coulomb-Kraft sein muss, also
\begin{equation}
                \frac{p^2}{mr} = \frac{e^2}{4\pi \epsilon_0} \frac{1}{r^2}
\end{equation}
($m$ Elektronenmasse, $e$ Elektronenladung, $\epsilon_0$ elektrische Feldkonstante),
erhalten wir zwei Bedingungen zwischen $p_n$ und $r_n$, sodass wir beide Gr\"o\ss en
als Funktion der Quantenzahl $n$ ausdr\"ucken k\"onnen:
\begin{equation}
         p_n = \frac{e^2 m}{4\pi \epsilon_0 \hbar} \frac{1}{n}  \hspace{1cm} \text{und} \hspace{1cm}
         r_n = \frac{4 \pi \epsilon_0 \hbar^2}{e^2 m} n^2\, . 
\end{equation}
Hier wurde angenommen, dass der Atomkern wesentlich schwerer ist als die Masse der
Elektronen. M\"ochte man die endliche Masse des Atomkerns ber\"ucksichtigen, muss man in
den folgenden Gleichungen die Elektronenmasse $m$ durch die reduzierte Masse
$\mu = m M/(m+M)$ ($M$ die Kernmasse) ersetzen.

Setzt man diese Beziehungen in die klassische Formel f\"ur die Energie ein,
\begin{equation}
                 E = \frac{p^2}{2m} - \frac{e^2}{4\pi \epsilon_0} \frac{1}{r} \, ,
\end{equation}
erhalten wir die nach der Quantisierungsbedingung erlaubten Energien  
in einem Wasserstoffatom:
\begin{equation}
                 E_n  = - \frac{e^4 m}{8 \epsilon_0^2 h^2} \frac{1}{n^2}  \, .
\end{equation}

Eine zweite wichtige Annahme von Bohr war, dass die beobachteten Frequenzen
der emittierten oder absorbierten Strahlung von einem Atom nicht den Energien der
m\"oglichen Zust\"ande entsprechen, sondern den \"Uberg\"angen zwischen zwei
m\"oglichen Energien in einem Atom. Das bedeutet, wenn ein Elektron von einem
erlaubten Zustand $E_n$ zu einem erlaubten Zustand $E_m$ springt, wird eine
Energie $\Delta E = E_m - E_n$ von dem Elektron in Form von Strahlung aufgenommen
bzw.\ abgegeben. Schlie\ss ich war die dritte Annahme von Bohr, dass diese
Energiedifferenz, die als Strahlung absorbiert bzw.\ emittiert wird, \"uber die
Formel von Planck und Einstein, $\Delta E = h \cdot \nu$, mit der
beobachteten Frequenz $\nu$ der Strahlung -- bzw.\ der beobachteten
Wellenl\"ande $\lambda = c/\nu$ -- in Beziehung steht. Damit waren nach dem
Bohr'schen Atommodell die folgenden Frequenzen als Absorptions- bzw.\
Emissionslinien erlaubt:
\begin{equation}
                 \nu(n,m) = \frac{1}{h} (E_n - E_m) = \frac{e^4 m}{8 \epsilon_0^2 h^3} 
                          \left( \frac{1}{n^2} - \frac{1}{m^2} \right)  \hspace{1cm} (m,n=1,2,3,...) \, .
\end{equation}

Sommerfeld konnte dieses Modell auf elliptische
Bahnkurven erweitern und damit das Verhalten von Wasserstoff in einem magnetischen
oder auch einem elektrischen Feld erkl\"aren. Doch schon f\"ur das Heliumatom schlugen
alle Erkl\"arungsans\"atze fehl. Und auch f\"ur die Atome mit mehreren Elektronen konnten
keine Bahnkurven f\"ur die Elektronen konstruiert werden. Es zeigte sich, dass die
Ans\"atze zur Konstruktion von erlaubten Bahnkurven f\"ur die Elektronen nur f\"ur integrable
Modelle funktionierten, bei denen die Bahnkurven periodisch sind. 

\subsection{Heisenbergs Umdeutung -- Juli 1925}

Wie schon erw\"ahnt, versucht Heisenberg, die neue Quantenmechanik ausschlie\ss lich
auf beobachtbare Gr\"o\ss en zu bauen. Darunter versteht er in erster Linie die Intensit\"aten
und die Frequenzen der Emissions- und Absorptionslinien. Er verbindet mit den klassischen
Variablen $x$ und $p$ (Ort und Impuls) neue Gr\"o\ss en $X$ und $P$, die direkt
mit diesen beobachtbaren Intensit\"aten und Frequenzen zusammenh\"angen sollen.  

Heisenberg m\"ochte in seiner Arbeit vom Juli 1925 zun\"achst eine vergleichsweise 
harmlose Frage kl\"aren: Wenn ein klassisches Objekt $x$ durch ein quantenmechanisches
Objekt $X$ ersetzt wird, durch welches quantenmechanische Objekt wird dann $x^2$
ersetzt? Zur Kl\"arung dieser Frage betrachtet er zun\"achst eine Bahnkurve $x(t)$ aus 
der klassischen Physik. Wir beschr\"anken uns hier auf periodische Bahnkurven (Heisenberg
betrachtet auch den aperiodischen Fall, allerdings ist dieser etwas problematisch), sodass
wir f\"ur die Zeitabh\"angigkeit dieser Kurve eine Fourier-Zerlegung ansetzen k\"onnen:
\begin{equation}
           x(t) = \sum_m x(m) {\rm e}^{2\pi {\rm i} m \nu t}  \, .
\end{equation}
(Die Summation erstreckt sich \"uber alle $m\in \mathbb{Z}$, wobei aus der Realit\"at von
$x(t)$ die Einschr\"ankung $x(-m)=x^*(m)$ folgt.)
$\nu$ ist die Grundfrequenz und $m \nu$ sind die (harmonischen) Oberfrequenzen. F\"ur
einen besseren  Vergleich zwischen der klassischen Mechanik und den \glqq umgedeuteten\grqq\
quantenmechanischen Gr\"o\ss en kennzeichnet Heisenberg eine solche Trajektorie noch
durch ihren Drehimpuls $J$, wobei er $J=n\hbar$ schreibt und den Index $n$ f\"ur eine
klassische Trajektorie noch kontinuierlich l\"asst. Damit erh\"alt er
\begin{equation}
           x(n;t) = \sum_m x(n,m) {\rm e}^{2\pi {\rm i} m \nu(n) t}  \, .
\end{equation}
Bildet man das Quadrat dieser Funktion, erh\"alt man:
\begin{equation}
           x(n;t)^2 = \sum_{m,k} x(n,m) x(n,k) {\rm e}^{2\pi {\rm i} (m+k) \nu(n) t} 
           = \sum_{m} \left( \sum_k x(n,m-k) x(n,k) \right) {\rm e}^{2\pi {\rm i} m \nu(n) t}  \, .
\end{equation}
Heisenberg leitet daraus ein klassisches Multiplikationsgesetz f\"ur
die Objekte $x(n,m)$ ab: $x^2(n,m)=\sum_k x(n,m-k) x(n,k)$. Dies entspricht einer
diskreten Faltung im zweiten Index. 

Heisenberg m\"ochte nun die klassischen Objekte $x(n,m)$ (die Amplituden) und $\nu(n)$ (die
Frequenzen) durch die quantenmechanisch beobachtbaren Gr\"o\ss en $X(n,m)$ und 
$\nu(n,m)$ ersetzen, die er mit den Intensit\"aten und Frequenzen zu einem \"Ubergang von einem 
Zustand $m$ in einen Zustand $n$ in Verbindung bringt. Die Frequenzen sollen dabei dem
Ritz'schen Kombinationsprinzip gen\"ugen, d.h., zu $\nu(n,m)$ soll es Frequenzen
$\nu(n)$ und $\nu(m)$ geben, sodass $\nu(n,m)=\nu(n)-\nu(m)$. F\"ur die Gr\"o\ss en $X(n,m)$ 
postuliert Heisenberg aus rein physikalischen \"Uberlegungen zur Kombination der \"Uberg\"ange 
zwischen den Zust\"anden die Regel:
\begin{equation}
\label{eq_1925_XX}
              X^2(n,m)= \sum_k X(n,k) X(k,m)  \, . 
\end{equation}
Er dr\"uckt diese Gleichung allerdings etwas anders aus (im Wesentlichen ist $m$ f\"ur ihn
die Differenz zwischen Anfangs- und Endzustand) und erkennt vermutlich nicht, dass
es sich hierbei um eine einfache Matrizenmultiplikation handelt. Matrizen geh\"orten damals
nicht zur Standardausbildung in der Physik. 

Heisenberg nutzt nun diese Multiplikationsregel, um beispielsweise die Korrekturen
zum anharmonischen Oszillator (also einem Oszillator, bei dem in
der ansonsten linearen Bewegungsgleichung ein Term proportional zu $x^2$ auftritt)
zu berechnen
und vergleicht diese mit bekannten Korrekturtermen.

Heisenberg formuliert f\"ur seine \glqq Umdeutungsvariablen\grqq\ $X$ das Programm:
Man bestimme zun\"achst L\"osungen der Bewegungsgleihungen f\"ur $X$, d.h., wenn die
klassischen Bewegungsgleichungen f\"ur die Trajektorien die Form
\begin{equation}
               \ddot{x} + f(x) = 0  
\end{equation} 
haben, dann sollen f\"ur $X$ dieselben Bewegungsgleichungen gelten. Er begr\"undet dies mit dem
Korrespondenzprinzip zwischen klassischen und quantenmechanischen Gr\"o\ss en.

Au\ss erdem soll $X$ eine Quantisierungsbedingung auferlegt werden, die -- ebenfalls
nach dem Korrespondenzprinzip -- das Analogon zur klassischen Quantisierungsbedingung
\begin{equation}
                   \oint p\,{\rm d}q=nh
\end{equation}
sein soll. Er ben\"otigt dazu das Produkt von $X$ mit $P$, den quantenmechanischen Gegenst\"ucken
zum Ort $x$ und zum Impuls $p$. Er stellt fest, dass das Produkt Gl.\ \ref{eq_1925_XX} f\"ur 
verschiedene Matrizen nicht kommutiert, also $XP$ nicht
dasselbe ist wie $PX$. Daher postuliert er f\"ur das Produkt eine Symmetriesierung:
$(XP+PX)/2$.
Er leitet dann speziell f\"ur den harmonischen Oszillator
(bei dem die klassischen Bahnkurven periodisch sind) eine Beziehung ab, die in
versteckter Form schon die kanonischen Vertauschungsrelationen enth\"alt. Klar
herausgearbeitet werden diese allerdings erst in der Arbeit von Born und Jordan.

\subsection{Born und Jordan -- September 1925}

Max Born hatte Mathematik und Physik studiert und war seit 1921 Professor f\"ur theoretische
Physik in G\"ottingen. Pascual Jordan war 1925 sein Assistent. Beide
kannten Matrizen und erkannten in den Multiplikationsgesetzen der quantenmechanischen
Objekte von Heisenberg die Matrixmultiplikation. Sie hatten die Ideen von
Heisenberg auch im \textit{statu nascendi} (wie sie selbst schreiben) kennengelernt.

Ihre Arbeit beginnt mit einer ausf\"uhrlichen Einf\"uhrung in die Theorie der
Matrizen. Sie formulieren die heute als Heisenberg'schen Bewegungsgleichungen
bekannten Gleichungen. Dazu erl\"autern sie zun\"achst, wie die Ableitung einer
Funktion einer Matrix $X$, also $f(X)$, nach der Matrix $X$ zu interpretieren ist. Dann
fordern Sie als Teil des sogenannten Korrespondenzprinzips, dass f\"ur eine
Hamilton-Matrix $H$, aufgefasst als Funktion von $X$ (dem Quantenanalogon zur kartesischen 
Koordinate) und $P$  (dem Quantenanalogon zum Impuls $p$), die Matrizen $X$ und $P$ 
die Gleichungen
\begin{equation}
       \dot{X} =  \frac{\partial H(X,P)}{\partial P} \hspace{1cm} {\rm und} \hspace{1cm} 
              \dot{P} = - \frac{\partial H(X,P)}{\partial X}  
\end{equation}
erf\"ullen sollen. Das sind genau die Bewegungsgleichungen, die wir heute als
Heisenberg'sche Bewegungsgleichungen bezeichnen.

Von besonderer Bedeutung ist, dass sie die
\glqq alte\grqq\ Quantisierungsbedingung
\begin{equation}
              \oint p \,{\rm d}q = n h
\end{equation}
durch eine neue, den Matrizen $P$ und $X$ angepasste Bedingung 
ersetzen, von der sie dann zeigen, dass sie \"aquivalent ist zu
\begin{equation}
              \sum_k \big( P(n,k) X(k,m) - X(n,k) P(k,m) \big) =  \frac{h}{2\pi {\rm i}} \, ,  
\end{equation}
also dem, was wir heute als die kanonischen Vertauschungsrelationen bezeichnen.
Heisenberg hatte diese Beziehung nur an einem Spezialfall (dem harmonischen Oszillator)
in versteckter Form erhalten.

Sie berechnen in ihrer Arbeit auch das Spektrum des Hamilton-Operators f\"ur
den harmonischen Oszillator. Au\ss erdem zeigen Sie, dass die Quadrate der
Matrixelemente $| X(n,m)|^2$ mit den Wahrscheinlichkeiten f\"ur einen durch
die Koordinate $x$ (einen Dipol) induzierten \"Ubergang
vom Energiezustand zu der Frequenz $\nu(m)$ zur Energie zu der Frequenz $\nu(m)$
zusammenh\"angt. 

Gelegentlich wurde spekuliert, weshalb Born und Jordan ihre Arbeit nicht zusammen
mit Heisenberg geschrieben haben. Hierzu sollte man wissen, dass Heisenberg Mitte Juni 
nach seiner R\"uckkehr aus Helgoland, wo er wichtige Rechnungen zu seinen
neuen Gr\"o\ss en vorgenommen hatte, ein Exemplar seiner Arbeit Born gegeben
und diesen gebeten hatte, die Arbeit einzureichen, sofern er den Inhalt als
interessant erachtete. Anschlie\ss end fuhr Heisenberg bis Ende September in Urlaub
und hatte mit Born nicht mehr gesprochen. Born hatte sich die Arbeit von Heisenberg zusammen
mit Jordan genauer angeschaut und diese dann auch eingereicht. Erst in diesem Zusammenhang
erkannten Born und Jordan, dass Heisenberg im Wesentliche eine Matrixmultiplikation
f\"ur seine Gr\"o\ss en formuliert hatte. Die Ende September eingereichte Arbeit von Born und
Jordan beruhte daher auf keinen weiteren Diskussionen mit Heisenberg und erschien deshalb
wohl ohne ihn als Koautor. Erst nach Heisenbergs R\"uckkehr arbeiteten alle drei an der
allgemeinen Formulierung zusammen und ver\"offentlichten diese Ergebnisse in der
\glqq Drei-M\"anner-Arbeit\grqq.

\subsection{Dirac -- November 1925}

Offenbar hatte Max Born ein Exemplar der Arbeit von Heisenberg auch an Ralph Fowler
geschickt, der gerade zum Fellow of the Royal Society ernannt worden war, und dieser
hatte es an Dirac weitergegeben, dessen Betreuer er war. Offenbar kannte Dirac die
sp\"atere Arbeit von Born und Jordan noch nicht, denn er bezieht sich ausschlie\ss lich
auf die Arbeit von Heisenberg, kommt aber im Wesentlichen zu den gleichen Ergebnissen
wie Born und Jordan. 

Sein Schwerpunkt war die nicht-kommutative Algebra, die
sich aus den Heisenberg'schen Multiplikationsregeln ergab. Anscheinend erkannte auch er
nicht, dass es sich um ein Matrixprodukt handelt, sondern bezeichnet die neue
Multiplikationsregel als \glqq Heisenberg-Produkt\grqq. Auch Dirac leitet zun\"achst wichtige
Regeln f\"ur das Matrixprodukt ab, insbesondere auch Ableitungsregeln. Eine wesentliche
Einsicht von Dirac ist, dass der Kommutator zweier Operatoren im Grenzfall gro\ss er 
Quantenzahlen und vergleichsweise kleiner Differenzen von Quantenzahlen bei 
atomaren \"Uberg\"angen in eine
klassische Poisson-Klammer \"ubergeht. Bei ihm findet man daher die Quantisierungsregel,
dass im Quantenformalismus die klassischen Poisson-Klammern durch den Kommutator
(multipliziert mit $h/2\pi {\rm i}$) zu ersetzen sind. 

Zun\"achst irritierend an seiner Arbeit ist, dass Dirac die Poisson-Klammer mit dem heutigen
Symbol f\"ur den Kommutator bezeichnet, d.h.\ $[x,y]$ ist die Poisson-Klammer f\"ur die beiden
Observablen $x$ und $y$. Der Kommutator wird immer explizit ausgeschrieben. Er formuliert
als Quantisierungsbedingung explizit die kanonischen Vertauschungsrelationen:
\begin{equation}
         q_r q_s - q_s q_r = 0   \hspace{1cm}  p_rp_s-p_sp_r=0 \hspace{1cm} 
         q_r p_s - p_s q_r = \delta_{rs}\, {\rm i} \frac{h}{2\pi}
\end{equation}
Dirac verifiziert, dass der Kommutator von zwei Gr\"o\ss en dieselben Beziehungen erf\"ullt
wie die Poisson-Klammer (insbesondere die Antisymmetrie und die Jacobi-Identit\"at) und
folgert daraus, dass f\"ur die quantisierten Gr\"o\ss en ganz allgemein dieselben Beziehungen
gelten m\"ussen wie im klassischen Fall. Schon in seiner Einleitung betont er, was erst in
der Arbeit abgeleitet wird, dass die Beziehungen zwischen beobachtbaren Gr\"o\ss en in der
klassischen Physik und in der Quantenphysik dieselben sind, und dass lediglich die
algebraischen Regeln -- das Multiplikationsgesetz -- f\"ur die quantenmechanischen
Gr\"o\ss en andere sind. 

\section{Die Wellenmechanik von Schr\"odinger 1926}

Wenn man heute die Schr\"odinger-Gleichung ableitet, geschieht das meist \"uber
die klassische Energiebeziehung
\begin{equation}
             E = \frac{p^2}{2m} + V(x) \, ,
\end{equation}
von der man verlangt, dass sie \"uber die de Broglie-Beziehungen $E=\hbar \omega$ und
$p= \hbar k$ zu einer Dispersionsrelation zwischen der Winkelfrequenz $\omega$ und
dem Wellenzahlvektor $k$ f\"ur ein Feld $\psi(x,t)$ wird, d.h., $\omega$ erh\"alt man durch
die Ableitung nach der Zeit $t$ und $k$ durch die Ableitung nach dem Ort $x$. 

Erwin Schr\"odinger hat \glqq seine\grqq\ Gleichung im Jahr 1926 jedoch nicht auf diese
Weise hergeleitet, sondern ist einen ganz anderen Weg gegangen. 
In diesem Jahr hat er mehrere Arbeiten zur Quantentheorie ver\"offentlicht, wobei zun\"achst
von besonderer Bedeutung eine Serie von vier Arbeiten (er nannte
sie \glqq Mitteilungen\grqq) mit dem Titel \glqq Quantisierung als Eigenwertproblem\grqq\ 
\cite{Schroedinger1bis4} ist. In diesen Arbeiten entwickelte er seine Idee von einer
Wellengleichung, die in ihrer zeitunabh\"angigen Form eine Eigenwertgleichung
f\"ur die m\"oglichen Energien eines Systems darstellt. In einer weiteren Arbeit, die am 18.\ M\"arz 1926
bei den Annalen der Physik einging, zeigte Schr\"odinger die \"Aquivalenz 
seiner Formulierung der Quantentheorie mit der Formulierung von Heisenberg.

Schr\"odinger war beeindruckt von der Arbeit von Louis de Broglie und der
Idee, dass jede Form von Materie nicht nur Teilcheneigenschaften haben soll
sondern auch Welleneigenschaften. Dies erinnerte ihn an einen Formalismus,
der urspr\"unglich aus der Optik stammt, der von Hamilton aber auf die Mechanik
\"ubertragen worden war und als Hamilton-Jacobi-Formalismus bekannt war.
Dieser Formalismus soll kurz angedeutet werden.

\subsection{Von Fermat und Huygens zu Hamilton und Jacobi} 

In der Optik beschreiben wir die Ausbreitung von Licht oft nach zwei Prinzipien:
Zum einen verwenden wir in der geometrischen Optik gerne das Fermat'sche Prinzip,
wonach sich Licht in einer optisch inhomogenen Umgebung zwischen zwei Punkten
so ausbreitet, dass die Zeit f\"ur die Zur\"ucklegung dieser Wegstrecke minimal wird.
Das Ergebnis dieses Extremalprinzips sind Trajektorien, die zwei Punkte verbinden.
Diese Trajektorien bilden die Grundlage der geometrischen Optik.

Zum anderen verwenden wir das Huygens'sche Prinzip, wonach jeder Punkt einer
Wellenfront zum Ausgangspunkt einer kreisf\"ormigen Elementarwelle wird, sodass die
Einh\"ullende all dieser Elementarwellen die neue Wellenfront bildet. 
Auch wenn wir heute das Huygens'sche Prinzip allgemein als ein Prinzip der Ausbreitung 
von Wellen interpretieren (Huygens stellte sich wohl eher Sto\ss prozesse in einem mit
winzigen Kugeln angef\"ullten Medium vor), spielt beim Huygens'schen Prinzip die
Wellenl\"ange keine Rolle; insofern steht das Huygens'sche Prinzip in seiner
urspr\"unglichen Form noch zwischen der Wellenoptik und der geometrischen Optik. 

\begin{figure}[htb]
\begin{picture}(200,140)(0,0)
\put(100,10){\makebox(0,0){$\bullet$}}
\qbezier(100,10)(50,70)(20,120)
\qbezier(100,10)(70,70)(50,130)
\qbezier(100,10)(85,70)(80,135)
\qbezier(100,10)(100,70)(110,135)
\qbezier(100,10)(115,70)(150,130)
\qbezier(100,10)(130,70)(180,120)
\thicklines
\qbezier(25,100)(95,152)(170,98)
\qbezier(37,85)(90,125)(154,85)
\qbezier(48,65)(95,100)(140,68)
\put(105,5){\makebox(0,0){$x_0$}}
\put(20,126){\makebox(0,0){$x_1$}}
\put(50,136){\makebox(0,0){$x_2$}}
\put(80,141){\makebox(0,0){$x_3$}}
\put(110,141){\makebox(0,0){$x_4$}}
\put(150,136){\makebox(0,0){$x_5$}}
\put(180,126){\makebox(0,0){$x_6$}}
\put(20,120){\makebox(0,0){$\bullet$}}
\put(50,130){\makebox(0,0){$\bullet$}}
\put(80,135){\makebox(0,0){$\bullet$}}
\put(110,135){\makebox(0,0){$\bullet$}}
\put(150,130){\makebox(0,0){$\bullet$}}
\put(180,120){\makebox(0,0){$\bullet$}}
\end{picture}
%
\hfill
\begin{picture}(210,140)(0,0)
\put(40,58.2){\makebox(0,0){$\bullet$}}
\put(50,60.2){\makebox(0,0){$\bullet$}}
\put(60,62.2){\makebox(0,0){$\bullet$}}
\put(70,63.4){\makebox(0,0){$\bullet$}}
\put(80,64.5){\makebox(0,0){$\bullet$}}
\put(90,65.0){\makebox(0,0){$\bullet$}}
\put(100,65.1){\makebox(0,0){$\bullet$}}
\put(110,65.0){\makebox(0,0){$\bullet$}}
\put(120,64.5){\makebox(0,0){$\bullet$}}
\put(130,63.4){\makebox(0,0){$\bullet$}}
\put(140,62.2){\makebox(0,0){$\bullet$}}
\put(150,60.2){\makebox(0,0){$\bullet$}}
\put(160,58.2){\makebox(0,0){$\bullet$}}
\put(40,58.2){\circle{30}}
\put(50,60.2){\circle{30}}
\put(60,62.2){\circle{30}}
\put(70,63.4){\circle{30}}
\put(80,64.5){\circle{30}}
\put(90,65){\circle{30}}
\put(100,65.1){\circle{30}}
\put(110,65){\circle{30}}
\put(120,64.5){\circle{30}}
\put(130,63.4){\circle{30}}
\put(140,62.2){\circle{30}}
\put(150,60.2){\circle{30}}
\put(160,58.2){\circle{30}}
\thicklines
\qbezier(10,50)(100,80)(190,50)
\qbezier(10,50.5)(100,80.5)(190,50.5)
\qbezier(10,66.5)(100,96.5)(190,66.5)
\end{picture}
\caption{\label{fig_HJ_1}%
Die Beziehung zwischen dem Fermat'schen Ausbreitungsprinzip und dem
Huygens'schen Ausbreitungsprinzip von Licht. (links) Nach dem Fermat'schen
Prinzip breiten sich von einem Punkt $x_0$ die Strahlen zu Punkten $x_i$ so aus,
dass die ben\"otigte Zeit minimal ist. Man kann nun jedem Punkt $x$ im Raum die Zeitdauer
$t$ der Trajektorie zuordnen, die das Licht von $x_0$ zu diesem Punkt $x$ ben\"otigt. Die
Linien konstanter Zeit bilden dann Fl\"achen, die von den Trajektorien senkrecht
durchsto\ss en werden. (rechts) Nach dem Huygens'schen Prinzip breiten sich von
jedem Punkt einer Wellenfront (das sind die Fl\"achen konstanter Zeit bei Fermat) 
kreisf\"ormige Elementarwellen aus, deren Einh\"ullende die neue Wellenfront bildet.}
\end{figure}
 
Ausgehend vom Fermat'schen Prinzip gelangt man zu den Huygens'schen Wellenfronten,
indem man ein Feld bestimmt, das zu jedem Raumpunkt $x$ die Zeit angibt, die das
Licht von einem festen Ausgangspunkt $x_0$ bis zu diesem Punkt $x$ ben\"otigt. Fl\"achen
konstanter Zeit bilden dann die Huygens'schen Wellenfronten. Da sich die Frequenz von Licht
w\"ahrend seiner Ausbreitung nicht \"andert (die Wellenl\"ange \"andert sich mit der optischen
Dichte eines Mediums nach der Gleichung $\nu \lambda(x) = c(x)$), 
sind Fl\"achen konstanter Zeit auch gleichzeitig Fl\"achen konstanter
Phase. Umgekehrt kann man
von dem Feld zu diesen Wellenfronten das Gradientenfeld bilden. Dieses definiert die Tangenten an die
Fermat'schen Trajektorien.

\subsection{Das Fermat'sche Prinzip}

Die Geschwindigkeit von Licht $c(x)$ an einem Punkt $x$ in einem Medium ist durch die optische 
Dichte bzw.\ den Brechungsindex $n(x)$ gegeben: $c(x) = c/n(x)$, wobei $c$ die Lichtgeschwindigkeit
im Vakuum ist. In optisch anisotropen Medien handelt es sich bei $n(x)$ um einen Tensor; diese
Problematik soll hier aber au\ss er Acht gelassen werden. Die Zeit $t(x)$, die das Licht von einem
Punkt $x_0$ zu einem Punkt $x$ entlang einer Trajektorie $\gamma=x(t)$ ben\"otigt, ist
\begin{equation}
                     T[\gamma](x,x_0) =  \int_{\gamma}  \frac{n(x')}{c}\,{\rm d}x' = \frac{1}{c} \int_{t(x_0)}^{t(x)} 
                       n(x(\tau))\, \left| \frac{{\rm d}x}{{\rm d}\tau}\right| \, {\rm d \tau} \, .  
\end{equation}
Hierbei wurde ausgenutzt, dass f\"ur gen\"ugend kleine Zeitspannen, in denen die
Geschwindigkeit als konstant angesehen werden kann, ${\rm d}t = {\rm d}s/v$ gilt. 

Nach dem Fermat'schen Prinzip ist die Aufgabe nun folgende: Gesucht ist die Trajektorie
$x(t)$, die das Funktional $T[x(t)](x,x_0)$ bei festgehaltenem Anfangspunkt $x_0$ und
festgehaltenem Endpunkt $x$ minimiert. Angenommen, diese Trajektorie ist gefunden --
dies ist die Trajektorie, die das Licht in der N\"aherung der geometrischen Optik 
tats\"achlich durchl\"auft --, dann k\"onnen wir jedem Punkt $x$ im Raum diese Zeit
$T(x,x_0)={\rm min}_{[x(t)]}T[x(t)](x,x_0)$ zuordnen. Die Fl\"achen $T(x,x_0)={\rm const}$
zu einem festgehaltenen Punkt $x_0$ bilden die Huygens'schen Wellenfronten. 

\subsection{Das Huygens'sche Prinzip}
   
Kennt man umgekehrt diese Fl\"achen konstanter Zeit (bzw.\ konstanter Phase), 
so kann man an jedem Punkt
die Fl\"achennormale bilden (genauer den Gradienten des Wirkungsfeldes, also $\nabla S(x)$)
und zu diesem Vektorfeld die Trajektorien konstruieren. Diese Trajektorien sind die
Fermat'schen Trajektorien. Sie durchsto\ss en die Huygens'schen Wellenfronten
orthogonal (siehe Abb.\ \ref{fig_HJ_1} (links)).

Das Huygens'sche Prinzip besch\"aftigt sich mit der Ausbreitung dieser Fl\"achen zu konstanten
Zeiten. Jeder Punkt einer solchen Fl\"ache bildet den Ausgangspunkt f\"ur eine Elementarwelle (siehe
Abb.\ \ref{fig_HJ_1} (rechts)). Die Einh\"ullende all dieser Elementarwelle (die sich lokal jeweils
mit der durch die optische Dichte des Mediums gegebenen Lichtgeschwindigkeit ausbreiten)
bildet dann die n\"achste Fl\"ache konstanter Zeit. Auf diese Weise erh\"alt man die Fl\"achen
konstanter Zeiten und kann aus den Gradienten zu diesen Fl\"achen wieder das Vektorfeld zu   
den Trajektorien bestimmen. 

\subsection{Die Hamilton-Jacobi-Gleichung}

Wir beginnen die analytische Mechanik heute gerne mit dem Lagrange-Formalismus,
aus dem die Bewegungsgleichungen -- die Euler-Lagrange-Gleichungen -- abgeleitet werden.
Dieser Formalismus hat zwei gro\ss e Vorteile: (1) das Hamilton'sche Prinizip, wonach
das Integral \"uber die Lagrange-Funktion als Funktional der Bahnkurven f\"ur die realisierte
Bahnkurve station\"ar wird, ist unabh\"angig vom gew\"ahlten Koordinatensystem. Daher
gilt es auch f\"ur verallgemeinerte Koordinaten, die den Zwangsbedingungen des Problems
angepasst sind. (2) Das Prinzip l\"asst sich auch relativistisch invariant formulieren (z.B.\
in einer Feldtheorie), wohingegen der Hamilton'sche Formalismus die Auszeichnung einer
Zeitkoordinate erfordert. 

Gegeben sei eine Lagrange-Funktion $L(q,\dot{q},t)$, die oft die Form
\begin{equation}
          L(q,\dot{q},t) = T(q(t),\dot{q}(t)) - V(q(t),t)
\end{equation}    
annimmt. $T$ ist die kinetische Energie ausgedr\"uckt in den verallgemeinerten Koordinaten
$q$ und $\dot{q}$ und meist eine quadratische Funktion von $\dot{q}$. $V$ ist die potenzielle
Energie, die oft nur von $q$ abh\"angt, aber eventuell noch explizit von der Zeit $t$ abh\"angen
kann. Wir definieren nun das Wirkungsfunktional:
\begin{equation}
\label{eq_HJ_S}
         S[q(t)](q_1,t_1,q_0,t_0) = \int_{t_0}^{t_1} L(q(t),\dot{q}(t), t) \, {\rm d} t  \, .
\end{equation}
Das Wirkungsfunktional ist ein Funktional von Wegen $q: [t_0,t_1] \rightarrow \mathbb{R}^n$,
die alle zum Zeitpunkt $t_0$ an einem Punkt $q_0=q(t_0)$ beginnen und zum Zeitpunkt
$t_1$ am Punkt $q_1=q(t_1)$ enden. Diese funktionale Abh\"angigkeit wird durch die eckige
Klammer $[q(t)]$ ausgedr\"uckt. Zur Ableitung der Euler-Lagrange-Gleichungen bestimmt man die 
erste Variation von $S[q(t)]$ 
unter den genannten Randbedingungen, Anfangs- und Endpunkt also festgehalten werden. 
    
Wir betrachten nun die Wirkung als eine Funktion des Orts $q$ und der Zeit $t$, bei dem eine 
Trajektorie endet, die bei $q_0$ zum Zeitpunkt $t_0$ beginnt. Bei dieser
Trajektorie soll es sich um eine L\"osung der Euler-Lagrange-Gleichung handelt. Das bedeutet:
\begin{equation}
       S(q,t;q_0,t_0) = {\rm min}_{[q(t)]} S[q(t)](q,t,q_0,t_0) 
                  = {\rm min}_{[q(t)]}  \int_{t_0}^{t_1} L(q(t),\dot{q}(t), t) \, {\rm d} t  \, .
\end{equation}
Wir setzen also in Gl.\ \ref{eq_HJ_S} eine L\"osung der Euler-Lagrange-Gleichung ein und
betrachten die Wirkung nicht mehr als Funktional von Trajektorien sondern als Funktion
der Endpunkte f\"ur die physikalisch realisierten Trajektorien. Die Ausgangsdaten, $q_0$ und $t_0$,
werden festgehalten, sodass die Wirkung zu einem Feld $S(q,t)$ wird. 

Statt nun dieses Feld dadurch zu berechnen, dass man f\"ur jeden Punkt $(q,t)$ die
L\"osung der Euler-Lagrange-Gleichung bestimmt, diese L\"osung in die
Lagrange-Funktion einsetzt und das Wirkungsintegral berechnet, kann man direkt eine
partielle Differentialgleichung f\"ur $S(q,t)$ aufstellen.
Man erh\"alt diese Differentialgleichung, indem man zun\"achst die Legendre-Transformierte
$H(q,p,t)$ von $L(q,\dot{q},t)$ bez\"uglich der Variablen $\dot{q}$ bestimmt, also
\begin{equation}
        H(q,p,t) = {\rm min}_{\dot{q}} (p \dot{q} - L(q,\dot{q},t)) \, .
\end{equation}
Dies ist die Hamilton-Funktion. Ersetzt man nun den Impuls $p$ durch den Gradienten
von $S$,
\begin{equation}
\label{eq_HJ_p}
        p = \nabla S(q,t) \, ,
\end{equation}
und ber\"ucksichtigt, dass die Zeitableitung der Wirkung gleich der Hamilton-Funktion
ist, erh\"alt man
\begin{equation}
        H\left( q, \frac{\partial S}{\partial q},t\right) = - \frac{\partial S}{\partial t} \, . 
\end{equation}
Dies ist die Hamilton-Jacobi-Gleichung. Wurde die Funktion $S(q,t)$ gefunden, 
erh\"alt man daraus eine gew\"ohnliche Differentialgleichung erster Ordnung
f\"ur die Impulse (Gl.\ \ref{eq_HJ_p}).

\subsection{Der Ansatz von Schr\"odinger}

In seiner ersten Mitteilung vom Januar 1926 l\"ost Schr\"odinger das Wasserstoffproblem
und zeigt, dass er mit seiner Formulierung, die sich aus einem Extremalprinzip ergibt,
dieselben Energiewerte erh\"alt wie sie aus der Bohr'schen Quantisierungsbedingung folgen.

Die Idee von Schr\"odinger ist folgende: Ausgangspunkt ist die Hamilton-Jacobi-Gleichung
(s.o.) f\"ur das Wasserstoffproblem. Das Feld $S(q,t)$ fasst er als (komplexe) Phase einer
Welle auf, wie es auch bei einer Kurzwellenasymptotik gemacht wird:
\begin{equation}
                 \psi(q,t) = \exp \frac{S(q,t)}{K} \, .
\end{equation} 
Die Konstante $K$ f\"uhrt er aus Dimensionsgr\"unden ein -- sie muss die Dimension einer
Wirkung haben -- und setzt sie sp\"ater
gleich $h/2\pi$, um die volle \"Ubereinstimmung mit den Bohr'schen Energiewerten zu
erhalten. Au\ss erdem nutzt er aus, dass $\frac{\partial S}{\partial t}=-E$ ist. Mit
\begin{equation} 
       \frac{\partial S}{\partial q} =  \frac{K}{\psi(q,t)} \frac{\partial \psi(q,t)}{\partial q} 
\end{equation}
ergibt sich f\"ur die Hamilton-Jacobi-Gleichung:
\begin{equation} 
       \left( \frac{\partial \psi}{\partial x} \right)^2 +  \left( \frac{\partial \psi}{\partial y} \right)^2 +
              \left( \frac{\partial \psi}{\partial z} \right)^2
               -  \frac{2m}{K^2} \left( E + \frac{e^2}{r} \right) \psi(q,t)^2 = 0   \, .
\end{equation}
Bisher hat Schr\"odinger nur eine Umformulierung der klassischen Hamilton-Jacobi-Gleichung
vorgenommen. Nun definiert Schr\"odinger eine \glqq Wirkung\grqq\ zu dieser Feldgleichung,
indem er einfach das r\"aumliche Integral bildet:
\begin{equation} 
   J =   \int \! {\rm d}x\, {\rm d}y\, {\rm d}z \left\{  
        \left( \frac{\partial \psi}{\partial x} \right)^2 +  \left( \frac{\partial \psi}{\partial y} \right)^2 +
              \left( \frac{\partial \psi}{\partial z} \right)^2
               -  \frac{2m}{K^2} \left( E + \frac{e^2}{r} \right) \psi(x,y,z,t)^2 \right\}  \, .
\end{equation}
Aus dieser Wirkung formuliert Schr\"odinger eine Quantisierungsbedingung, die darin besteht, dass
die erste Variation verschwinden soll. Dies f\"uhrt auf die Gleichung
\begin{equation}
        \Delta \psi + \frac{2m}{K^2} \left( E + \frac{e^2}{r} \right) \psi = 0
\end{equation} 
bzw., mit $K=h/2\pi$:
\begin{equation}
     -  \frac{\hbar^2}{2m}  \Delta \psi(q,t)  - \frac{e^2}{r} \psi(q,t) =  E \psi (q,t) \, .
\end{equation} 
Dies ist die Schr\"odinger-Gleichung in ihrer heutigen Form. 
Schr\"odinger erkannte, dass es sich hier um eine Eigenwertgleichung handelt, daher
postuliert er als Quantisierungsbedingung die L\"osungen dieser Eigenwertgleichung. 
Er zeigt, dass es f\"ur positive Energien und physikalisch sinnvolle Randbedingungen 
immer L\"osungen gibt, dass die negativen
Energien aber diskret sind und mit den Bohr'schen Energien \"ubereinstimmen. 

Der Hamilton-Jacobi-Formalismus enth\"alt keine Wellenl\"ange, ebenso wie
auch das Huygens'sche Prinzip f\"ur die Optik unabh\"angig von einer Wellenl\"ange
ist. Schr\"odinger war aber bekannt, dass sich das Huygens'sche Prinzip bzw.\ die
Eikonalgleichung (die \glqq Hamilton-Jacobi-Gleichung\grqq, die sich aus dem
Fermat'schen Extremalprinzip ergibt) als
der f\"uhrende Term einer Kurz\-wellen\-asymp\-totik, also einer N\"aherung f\"ur
sehr kurze Wellenl\"angen (im Vergleich zu den Abmessungen von optischen
Hindernissen oder Inhomogenit\"aten), aus den Maxwell-Gleichungen ergibt.

W\"ahrend Schr\"odinger in seiner ersten Mitteilung vom Januar 1925 \glqq seine\grqq\
Gleichung aus einem Variationsprinzip ableitet, dessen Wirkungsfunktional er mehr oder
wenige willk\"urlich als
Interal \"uber die Hamilton-Jacobi-Gleichung schreibt, versucht er in seiner
zweiten Mitteilung vom 23.\ Februar 1926 seine Feldgleichung aus einer Wellengleichung 
(mit einem d'Alembert-Operator) abzuleiten. Erst in seiner vierten Mitteilung formuliert
Schr\"odinger die zeitabh\"angige Version in der Form, wie wir sie heute kennen.

In seiner zweiten Mitteilung geht er auch ausf\"uhrlicher
auf die Beziehungen zur Optik -- die Beziehung zwischen Fermat'schem und Huygens'schem
Prinzip -- ein. Au\ss erdem l\"ost er auch das Eigenwertproblem f\"ur den harmonischen
Oszillator sowie f\"ur den rotierenden K\"orper. 

\section{Die \"Aquivalenz zwischen Schr\"odinger und Heisenberg}
 
Die L\"osungen \glqq seiner\grqq\ Gleichung waren Wellenfunktionen $\psi_n(x)$
zu den m\"oglichen Energien $E_n$. Die wesentliche Beziehung zwischen
dem Heisenberg'schen Formalismus, bei dem jeder klassischen Gr\"o\ss e (Ort, Impuls,
Energie, Drehimpuls, etc.) eine Matrix -- ein mathematisches Objekt mit zwei Indizes -- 
zugeordnet wurde,\index{Heisenberg-Bild}\index{Matrizenmechanik}
und seinem Formalismus bestand in der Beziehung
\begin{equation}
         F(m,n) = \int_{V} \psi_m(x) f(x,-{\rm i}\hbar \partial/\partial x) \psi_n(x)\, {\rm d}V \, .
\end{equation}
Hierbei ist $f(x,p)$ eine klassische beobachtbare Gr\"o\ss e, ausgedr\"uckt als Funktion
des Orts $x$ und des Impulses $p$. Den Impuls $p$ ersetzt Schr\"odinger durch den
Operator $-{\rm i}\hbar \partial/\partial x$, wobei er die Beziehung von de Broglie zwischen
dem Impuls und der Wellenl\"ange ausnutzt. Wenn, nach de Broglie, ein Elektron mit einem
scharfen Impuls $p$ durch eine Welle mit einer scharfen Wellenl\"ange $\lambda=p/h$
und damit durch
\begin{equation} 
           \psi(x)=\exp  \left(  {\rm i}\frac{px}{h} \right)
\end{equation}  
zu beschreiben ist, dann erh\"alt man den Impuls durch die Ableitung nach $x$, multipliziert
mit $h$ (und der imagin\"aren Einheit i). Schr\"odinger
diskutiert ausf\"uhrlich, weshalb man die Funktion $f$ in ihren Argumenten $x$ und $p$
noch symmetrisieren muss, da die Multiplikation mit $x$ und die Ableitung nach $x$ nicht
vertauschen. Die
Gr\"o\ss e $F(m,n)$ entspricht dann der Heisenberg'schen Matrix, durch die man die
klassische Gr\"o\ss e in der Quantentheorie zu ersetzen hat.

Zum Zeitpunkt dieser Publikation kannte Schr\"odinger alle drei Arbeiten von Heisenberg, Born
und Jordan, und er bezieht sich auch darauf. Er
hatte erkannt, dass das allgemeine Problem -- bestimme das Spektrum
des Energieoperators $H= \frac{1}{2m} P^2 + V(Q)$ f\"ur Operatoren $P$ und $Q$, welche
den kanonischen Vertauschungsrelationen $[Q,P]=-{\rm i}\hbar$ gen\"ugen -- \"aquivalent ist zu
dem Problem, die Eigenwertgleichung f\"ur den Operator $H=-\frac{\hbar^2}{2m} \Delta
+V(x)$ zu l\"osen. Insbesondere erkannte er, dass die Operatoren \glqq Multiplikation mit $x$\grqq\
und $-{\rm i}\hbar \partial/\partial x$ die von Born und Jordan gefundenen kanonischen
Vertauschungsregeln erf\"ullen. 


\begin{thebibliography}{99}
\bibitem{Bohr1} Bohr, N.; \textit{On the Constitution of Atoms and Molecules};
         Philosophical Magazine {\bf 26},1 (1913).           
\bibitem{BornJordan} Born, M., Jordan, P.; \textit{Zur Quantenmechanik};
         Zeitschrift f\"ur Physik Bd.\ 34 (1925) 858--888.           
\bibitem{DreiMann} Born, M., Heisenberg, W.,
       Jordan, P.; {\em Zur Quantenmechanik II}; Zeitschrift f.\ Physik 35
       (1925) 557--615.       
\bibitem{Dirac} Dirac, P.; \textit{The fundamental equations of quantum mechanics}; Proceedings Royal
            Society, Band 9, 1925, S.\ 642--653.
\bibitem{Heisenberg} Heisenberg, W.; \textit{\"Uber quantentheoretische Umdeutung kinematischer
            und mechanischer Beziehungen}; Zeitschrift f\"ur Physik Bd.\ 33 (1925) 879--893.     
\bibitem{Pauli} Pauli, W.; \textit{\"Uber das Wasserstoffspektrum vom Standpunkt der neuen
             Quantenmechanik}; Zeitschrift f\"ur Physik, Band 36 (1926) 336--363.            
\bibitem{Schroedinger1bis4} Schr\"odinger, E.; \textit{Quantisierung als Eigenwertproblem};
           Annalen der Physik 1926, 4.\ Folge; Erste Mitteilung, Bd.\ 79, S.\ 361--376 (eingegangen 27.\ Januar 1926),
           Zweite Mitteilung, Bd.\ 79, S.\ 489--527 (eingegangen 23.\ Februar 1926), Dritte Mitteilung,
           Bd.\ 80, S.\ 437--490 (eingegangen 10.\ Mai 1926), Vierte Mitteilung, Bd.\ 81, S.\ 109--139 (eingegangen
           21.\ Juni 1926).
\bibitem{Schroedinger26} Schr\"odinger, E.; \textit{\"Uber das Verh\"altnis der Heisenberg-Born-Jordanschen
           Quantenmechanik zu der meinen}; Annalen der Physik 79 (1926) 734--756 (eingegangen 18.\ M\"arz 1925).              
\end{thebibliography}

\end{document}

